% !TEX program = xelatex
% Basic Arabic document with polyglossia + bidi + Amiri font.
% This is the minimal working setup for an Arabic XeLaTeX document.
\documentclass[12pt, a4paper]{article}

\usepackage{fontspec}
\usepackage{polyglossia}
\setmainlanguage[locale=mashriq]{arabic}
\setotherlanguage{english}

\newfontfamily\arabicfont[Script=Arabic]{Amiri}

\begin{document}

\title{مقدّمة في التنضيد العربي}
\author{د. نبراس أبو الذهب}
\date{2025}
\maketitle

\section{التمهيد}

اللغةُ العربيةُ من أكثرِ لغاتِ العالمِ انتشاراً، وهي لغةُ القرآنِ الكريم.
يَحتاجُ تنضيدُ النصوصِ العربيةِ في \LR{LaTeX} إلى إعدادٍ خاصٍّ بسببِ
الكتابةِ من اليمينِ إلى اليسارِ واتصالِ الحروفِ فيما بينها.

تُتيحُ حزمةُ \LR{polyglossia} معَ \LR{bidi} دعماً كاملاً للغةِ العربيةِ
في محركِ \LR{XeLaTeX}. هذا المثالُ يُوضّحُ الإعدادَ الأدنىَ المطلوبَ
لبدءِ كتابةِ مستندٍ عربيّ.

\end{document}
